\section{The spherical cosine and sine laws}
\label{sec:core-laws}

\subsection{The side cosine laws}

\begin{theorem}[Spherical law of cosines for sides]
\label{thm:side-cosine}
For a proper spherical triangle,
\begin{align}
 \cos a&=\cos b\cos c+\sin b\sin c\cos A,                 \tag{4.1}\\
 \cos b&=\cos c\cos a+\sin c\sin a\cos B,                 \tag{4.2}\\
 \cos c&=\cos a\cos b+\sin a\sin b\cos C.                 \tag{4.3}
\end{align}
Equivalently, without division,
\begin{align}
 s_bs_c c_A&=x-yz,&
 s_cs_a c_B&=y-zx,&
 s_as_b c_C&=z-xy.                                        \tag{4.4}
\end{align}
\end{theorem}

\begin{proof}
At \(A\), use the tangent vectors (3.7)--(3.8):
\[
 \cos A=t_{AB}\dotp t_{AC}
 =\frac{(v-zu)\dotp(w-yu)}{s_cs_b}
 =\frac{x-yz}{s_bs_c}.
\]
Multiplication by \(s_bs_c\) gives
\(x=yz+s_bs_c\cos A\), which is (4.1).  Cyclic relabeling gives
(4.2)--(4.3).  Every step before the final presentation is a polynomial dot-
product expansion; division uses the supplied positive side-sine witnesses.
\end{proof}

For a sphere of radius \(R>0\), the dimensionless central angles are
\(a/R,b/R,c/R\).  Thus (4.1) becomes
\[
 \cos\frac aR=\cos\frac bR\cos\frac cR+
 \sin\frac bR\sin\frac cR\cos A.                           \tag{4.5}
\]
The proof is unchanged after normalizing the radius vector by \(R\).

\begin{corollary}[Rearranged cosine formulas]
\label{cor:rearranged-cosine}
With the usual positive denominator witnesses,
\begin{align}
 \cos A&=\frac{\cos a-\cos b\cos c}{\sin b\sin c},         \tag{4.6}\\
 \cos a&=\frac{\cos A+\cos B\cos C}{\sin B\sin C}.         \tag{4.7}
\end{align}
The remaining formulas are cyclic.
\end{corollary}

The second formula is proved below as the dual cosine law.  Constructively,
(4.6) is shorthand for (4.4) plus the apartness of \(\sin b\sin c\).

\subsection{The spherical sine law}

\begin{theorem}[Spherical law of sines]
\label{thm:sine-law}
For a proper spherical triangle,
\begin{equation}
 \frac{\sin A}{\sin a}
 =\frac{\sin B}{\sin b}
 =\frac{\sin C}{\sin c}
 =\frac{\tau}{\sin a\sin b\sin c}.                         \tag{4.8}
\end{equation}
Its division-free content is
\begin{equation}
 \sin A\sin b=\sin B\sin a,
 \quad
 \sin B\sin c=\sin C\sin b,
 \quad
 \sin C\sin a=\sin A\sin c.                              \tag{4.9}
\end{equation}
\end{theorem}

\begin{proof}
Equation (3.10) gives
\[
 \sin A=\frac{\tau}{s_bs_c}.
\]
Dividing by \(s_a=\sin a\) yields the last expression in (4.8).
The two cyclic formulas have the same numerator and denominator.  For the
strictly algebraic version, cross multiply; for example
\[
 \sin A\sin b=\frac{\tau}{s_bs_c}s_b=\frac\tau{s_c}
 =\frac{\tau}{s_cs_a}s_a=\sin B\sin a.
\]
All canceled quantities have positive apartness witnesses.
\end{proof}

\begin{corollary}[Scalar-triple-product and volume forms]
\label{cor:sine-volume}
If \(V\) is the Euclidean volume of the tetrahedron with vertices
\(0,u,v,w\), then \(6V=\tau\), and
\begin{equation}
 \frac{\sin A}{\sin a}
 =\frac{6V}{\sin a\sin b\sin c}.                           \tag{4.10}
\end{equation}
This is the vector/triple-product form commonly used to prove the spherical
sine law.
\end{corollary}

\subsection{The angle cosine laws}

\begin{theorem}[Supplemental spherical law of cosines]
\label{thm:angle-cosine}
The angles and opposite sides satisfy
\begin{align}
 \cos A&=-\cos B\cos C+\sin B\sin C\cos a,                \tag{4.11}\\
 \cos B&=-\cos C\cos A+\sin C\sin A\cos b,                \tag{4.12}\\
 \cos C&=-\cos A\cos B+\sin A\sin B\cos c.                \tag{4.13}
\end{align}
Equivalently,
\begin{equation}
 \cos a=\frac{\cos A+\cos B\cos C}{\sin B\sin C},        \tag{4.14}
\end{equation}
with cyclic companions and explicit denominator witnesses.
\end{theorem}

\begin{proof}
We give a direct finite-algebra proof; Section~\ref{sec:polar} gives the more
conceptual proof by polar duality.  Substitute (3.11)--(3.12) into the
right-hand side of (4.11).  After multiplication by the positive quantity
\(s_a^2s_bs_c\), the desired equality becomes
\begin{equation}
 -(y-zx)(z-xy)+x\tau^2=(1-x^2)(x-yz).                      \tag{4.15}
\end{equation}
Insert the Gram identity (3.5); expansion reduces both sides to the same
integer polynomial.  Dividing by \(s_a^2s_bs_c\) gives (4.11).  The other
identities are cyclic.  Solving (4.11) for \(\cos a\) gives (4.14).
\end{proof}

\subsection{The complete core as a finite ideal}

The primary spherical theory can therefore be generated by the polynomial
relations
\begin{align*}
&s_a^2+x^2=1,\quad s_b^2+y^2=1,\quad s_c^2+z^2=1,\\
&\tau^2=1+2xyz-x^2-y^2-z^2,\\
&s_bs_c c_A=x-yz,\quad s_cs_a c_B=y-zx,\quad
 s_as_b c_C=z-xy,\\
&s_bs_c s_A=\tau,\quad s_cs_a s_B=\tau,\quad
 s_as_b s_C=\tau.                                          \tag{4.16}
\end{align*}
Every formula in this section is a consequence of (4.16) by polynomial
normalization and guarded cancellation.  This finite presentation is the core
object to formalize in Lean and in \(\ipra\).
